\section{Cipher Vigenère}

Cipher Vigenère adalah cipher polyalphabetic yang menggunakan kata kunci berulang. Setiap huruf plaintext digeser sesuai huruf kunci yang bersesuaian. Dengan kunci ``KEY'', huruf pertama digeser K (10), huruf kedua E (4), huruf ketiga Y (24), lalu berulang.

\begin{table}[!htbp]
\centering
\caption{Contoh Vigenère: plaintext ``HELLO'', kunci ``KEY''}
\label{tab:vigenere}
\begin{tabular}{@{}cccccc@{}}
\toprule
Plaintext & H & E & L & L & O \\
Kunci & K & E & Y & K & E \\
Ciphertext & R & I & J & V & S \\
\bottomrule
\end{tabular}
\end{table}
Rumus: $C_i = (P_i + K_i) \bmod 26$, dengan $P_i$, $K_i$ adalah posisi huruf (A=0).

\begin{lstlisting}[style=matlab,caption={Vigenère Cipher dalam MATLAB}]
function c = vigenere_encrypt(plain, key)
    plain = upper(plain); key = upper(key);
    c = '';
    for i = 1:length(plain)
        if plain(i) >= 'A' && plain(i) <= 'Z'
            ki = mod(i-1, length(key)) + 1;
            p = double(plain(i)) - 65;
            k = double(key(ki)) - 65;
            c = [c char(mod(p+k, 26) + 65)];
        else, c = [c plain(i)]; end
    end
end
\end{lstlisting}

Lebih kuat dari Caesar karena satu huruf plaintext dapat dipetakan ke banyak huruf ciphertext. Namun rentan terhadap \textbf{Kasiski examination} jika panjang kunci dapat diperkirakan \cite{schneier1996,trappe2006}.
